\documentclass[11pt]{article}
\usepackage[letterpaper,margin=.87in]{geometry}
\usepackage[T1]{fontenc}
\usepackage{lmodern,microtype,amsmath,amssymb,amsthm,booktabs}
\usepackage[hidelinks]{hyperref}
\newtheorem{theorem}{Theorem}
\setlength{\parindent}{0pt}
\setlength{\parskip}{.45em}
\raggedbottom
\title{\textbf{Three-color classification\\for shifted Schur equations}}
\author{1200}
\date{6 November 2026}
\begin{document}
\maketitle
\begin{abstract}
Assign a nonnegative shift to each of three colors and forbid monochromatic
solutions of the corresponding shifted Schur equation. We give one explicit
periodic coloring for every nonconstant triple of shifts. On the diagonal,
Schaal's formula gives the exact finite threshold.
\end{abstract}

\section{Statement}
Let $R_3(c_0,c_1,c_2)$ be the least $N\geq1$ such that every map
$[1,N]\to\{0,1,2\}$ has $x,y,z$ of a color $i$ with $x+y+c_i=z$.
Set $R_3=\infty$ if there is no such $N$. We allow $x=y$ and empty color classes.

\begin{theorem}
For nonnegative integer shifts,
\[
 R_3(c_0,c_1,c_2)=
 \begin{cases}
  13c+14,& c_0=c_1=c_2=c,\\
  \infty,& \text{otherwise}.
 \end{cases}
\]
In the nonconstant case there is an avoiding coloring of $\mathbb Z$ of period
$2^{s+1}$, where $s$ is the least $2$-adic valuation of a nonzero pairwise
difference of the shifts.
\end{theorem}

\section{A periodic coloring for unequal shifts}
Put $m=2^s$. The three shifts agree modulo $m$, and at least two differ
modulo $2m$. Exactly two therefore have the same residue $a$ modulo $2m$;
the remaining residue is $a+m$ modulo $2m$. Relabel colors so that the
first two agree. Under the translation $u=x+a-1$, $v=y+a-1$, $w=z+a-1$,
the equation becomes $u+v+(c_i-a+1)=w$. Thus the new shifts modulo $2m$
are $(1,1,1+m)$.

Write each representative as $hm+r$, with $h\in\{0,1\}$ and $0\leq r<m$.
Let $q\in\{1,-1\}$ represent $m$ modulo $3$. Define the color of a representative by
\[
 \chi(hm+r)=
 \begin{cases}r\bmod 3,&h=0,\\ (1+q-r)\bmod 3,&h=1.\end{cases}
\]
For a monochromatic solution of color $i$, set $H=h_u+h_v-h_w$ and
$\delta=\mathbf 1_{\{i=2\}}$. The residue equation gives
\[
 r_u+r_v+1-r_w=jm,\qquad j\equiv H+\delta\pmod 2.
\]
The left side lies in $[2-m,2m-1]$, so $j$ is $0$ or $1$, including when $m=1$.
A representative of color $i$ has $r\equiv i+(1+q-2i)h\pmod3$.
Consequently a solution would make $i+1+(1+q-2i)H-jq$ vanish modulo $3$.
The following table evaluates its residue; $j$ is the representative of
$H+\delta$ in $\{0,1\}$.
\[
\begin{array}{cc|rrrr}
 q&i&H=-1&H=0&H=1&H=2\\ \hline
 1&0&1&1&2&2\\
 1&1&1&2&1&2\\
 1&2&2&2&1&1\\
 -1&0&2&1&2&1\\
 -1&1&2&2&1&1\\
 -1&2&1&1&2&2
\end{array}
\]
No entry is zero. Pull the coloring back along $x\mapsto (x+a-1)\bmod 2m$
and undo the relabeling. Any integer solution would yield a modular solution,
so the coloring also works on every positive interval.

\section{The diagonal value}
The diagonal formula is due to Schaal~\cite{schaal}. Here is the scaling
argument of Adhikari et al.~\cite[Proposition 4.9 and Corollary 4.10]{adhikari}.
Set $b=c+1$ and translate each variable by $c$. The equation becomes
ordinary addition on $[b,N+c]$. For $N=13c+14$ this interval contains
$b,2b,\ldots,14b$. Every three-coloring of $[1,14]$ has a monochromatic
sum triple, by the classical three-color Schur value. Pulling the coloring
on the multiples of $b$ back to $[1,14]$ proves the upper bound.

For $N=13c+13$ the translated interval is $[b,14b-1]$. Color a translated
integer $y$ by the set containing $\lfloor y/b\rfloor$, using
\[
 \{1,4,10,13\},\quad \{2,3,11,12\},\quad \{5,6,7,8,9\}.
\]
For any one of these sets $A$ and $u,v\in A$, neither $u+v$ nor $u+v+1$
belongs to $A$. These are the only two possible block indices of the sum
of integers with block indices $u,v$. This proves the lower bound.

\section{Scope and reproducibility}
Adak et al.~\cite{adak} treat two colors with two assigned shifts. Their
classification does not cover the three-color problem. Schaal's formula
is the known diagonal input. A short companion script checks the residue
table, the carry condition and the Schur endpoint $14$. The proof is a
paper proof. It does not rely on the archived Lean prototype.

\textbf{AI contribution.} GLM5.3F (source project credit supplied by the
author). The periodic argument was simplified with GPT Sol Max assistance.

\begin{thebibliography}{9}
\bibitem{schaal} D. Schaal, A family of 3-color Rado numbers,
\emph{Congr. Numer.} \textbf{111} (1995), 150--160.
\bibitem{adhikari} S. D. Adhikari, L. Boza, S. Eliahou, J. M. Marín,
M. P. Revuelta and M. I. Sanz, On the $n$-color Rado number for the equation
$x_1+\cdots+x_k+c=x_{k+1}$, \emph{Math. Comp.} \textbf{85} (2016),
2047--2064. \href{https://doi.org/10.1090/mcom3034}{doi:10.1090/mcom3034}.
\bibitem{adak} R. Adak, Y. Bakshi, L. S. Chandran and S. G. Nanoti, Off-diagonal Rado numbers for
$x+y+c=z$ and $x+y+k=z$, \href{https://arxiv.org/abs/2603.28216}{arXiv:2603.28216} (2026).
\end{thebibliography}
\end{document}
